\tzpanchor{1077}
\tzpentryheader{1077}{MEASUREMENT PHASE INTEGRATION}

\begin{tzpsnapshot}
\tzpsnaprow{TZPID}{\tzpcode{ID1077}}
\tzpsnaprow{UUID}{\tzpcode{9fb3fb0b-3234-50ba-a0e2-76bf2a8a1553}}
\tzpsnaprow{Type}{Transfer Statement}
\tzpsnaprow{Scale}{transfer}
\tzpsnaprow{LOC Classification}{Working routing: QA641 manifolds and topology; QC174.17 mathematical physics}
\tzpsnaprow{arXiv Classification}{Primary: math-ph; Cross-list: gr-qc, math.DG}
\end{tzpsnapshot}

\tzpsection{Canonical Statement}
M = sum\_i |irangle langle i| otimes |irangle\_\{cl\} implies text\{Holevo Readout\}

\[
M = \sum_{\mathrm{i}} \lvert i \rangle \langle i| \otimes \lvert i \rangle_{cl} \implies \text{Holevo Readout}
\]
\[
mathcal{J}_{\mathrm{out}} = \mathcal{T}_{\mathrm{TZP}}\!\left(\mathcal{K}_{\mathrm{TZP}}\right)
\]

\noindent
\begin{minipage}[t]{0.48\linewidth}
\tzpsection{Wolfram Form}
\begingroup
\small
\[
\begin{aligned}
\mathfrak{W}_{\mathrm{TZP}}\!\left(\mathcal{K}_{1077}\right)
&\longmapsto
\left\langle \Phi_{\tau},\; \Omega^{\mathbb{T}},\; \Psi_{\mathrm{T}} \right\rangle
\end{aligned}
\]
\[
\mathcal{K}_{1077}
:=
\operatorname{HoldForm}\!\left(\mathcal{T}_{1077}\right)
\]
\[
\begin{aligned}
\operatorname{Admissible}_{\mathrm{TZP}}\!\left(\mathcal{K}_{1077}\right)
&\Longleftrightarrow
\mathfrak{W}_{\mathrm{TZP}}\!\left(\mathcal{K}_{1077}\right)\not\equiv\varnothing
\end{aligned}
\]
\textbf{Wolfram register:} canonical symbol lift\\
\textbf{Title lift:} MEASUREMENT PHASE INTEGRATION\\
\textbf{Symbol key:} \(\Phi_{\tau}\): phase-time symmetry\\ \(\Omega^{\mathbb{T}}\): Trawinistic boundary curvature\\ \(\Psi_{\mathrm{T}}\): Trawinistic wavefunction
\par
\endgroup
\end{minipage}\hfill
\begin{minipage}[t]{0.48\linewidth}
\tzpsection{TRAWIN Normalization}
\[
\tzpnormalizewrap{M = \sum_{\mathrm{i}} \lvert i \rangle \langle i| \otimes \lvert i \rangle_{cl} \implies \text{Holevo Readout},\,mathcal{J}_{\mathrm{out}} = \mathcal{T}_{\mathrm{TZP}}\!\left(\mathcal{K}_{\mathrm{TZP}}\right),\,mathcal{A}_{\mathrm{adm}}\!\left(\mathcal{J}_{\mathrm{out}}\right)=1}
\]

\small
\textbf{T}: Aggregate Relation is localized within the category theory and proof structure arc of the directory.\\
\textbf{R}: Support Relation preserves the operative relation that measurement phase integration requires.\\
\textbf{A}: Closure Relation fixes the admissible closure condition for the entry.\\
\textbf{W}: MEASUREMENT PHASE INTEGRATION is rendered as a reusable operator-level statement.\\
\textbf{I}: The informational content of measurement phase integration is retained rather than flattened.\\
\textbf{N}: MEASUREMENT PHASE INTEGRATION closes as a distinct normalized TRAWIN object.
\normalsize
\end{minipage}

\tzpsection{Derivable Consequences}
The entry now reads through actual sourced equations, so its local arc is carried by explicit mathematical relations rather than a uniform quaternary certificate record shell.
\clearpage

\tzpentryheader{1077}{Oxford Dictionary of the Queens, NY English}

\begingroup\small
\tzpsection{Definition}
MEASUREMENT PHASE INTEGRATION is a registry definition in Category Theory and Proof Structure with secondary pressure from Signal.

% BEGIN TZP CONTEXTUAL MEANING
\tzpsection{Contextual Meaning}
\begingroup\footnotesize\setlength{\emergencystretch}{3em}\sloppy
\textbf{Formal statement:} M = sum sub i |irangle langle i| otimes |irangle sub cl implies text Holevo Readout. \textbf{Contextual meaning:} The entry reads MEASUREMENT PHASE INTEGRATION as a controlled technical headword whose equation layer, proof route, and registry role are interpreted together. \textbf{Derivable consequences:} The entry now reads through actual sourced equations, so its local arc is carried by explicit mathematical relations rather than a uniform quaternary certificate record shell. \textbf{Lemma support:} Aggregate Relation; Support Relation; Closure Relation.\par
\endgroup
% END TZP CONTEXTUAL MEANING

\tzpsection{Technical Interpretation}
Technically, MEASUREMENT PHASE INTEGRATION is read through the local equation chain and the TRAWIN normalization recorded on the entry page.

\tzpsection{Equation Grounding}
Equation anchors: M = sum sub i lvert i i| otimes lvert i sub cl implies Holevo Readout; mathcal J sub out = T sub TZP ( K sub TZP ); and mathcal A sub adm ( J sub out )=1. Proof route: show that the stated equations form a coherent progression for the entry and remain compatible under the TRAWIN ordering. Formalization route: prove that the final closure relation follows from the preceding entry-specific equations.

\tzpsection{Interpretive Role}
The entry now reads through actual sourced equations, so its local arc is carried by explicit mathematical relations rather than a uniform quaternary certificate record shell.
\endgroup

\tzpsection{TZP Thesaurus}
\begin{enumerate}[leftmargin=1.6em,itemsep=6pt,topsep=4pt]
\item \textbf{Measurement Cycle Angle Integration}: quantum-state language for amplitudes, measurement, basis structure, or weak coupling.
\item \textbf{Cycle Angle Quantum-State Analogue}: quantum-state language for amplitudes, measurement, basis structure, or weak coupling.
\item \textbf{Cycle Angle Terminology}: entry-specific alternate wording for indexing, related literature, and local formal context.
\end{enumerate}

\tzpsection{Encyclopedia Mathematica}
\begin{samepage}
\noindent\begin{minipage}[t]{0.45\linewidth}
\vspace{0pt}
\raggedright
\scriptsize\textbf{Image reading.} The rendered manifold at right is the per-ID light plate for MEASUREMENT PHASE INTEGRATION. It should be read as the visual companion to the entry's equation layer, not as a repeated template diagram.\par
\end{minipage}\hfill
\begin{minipage}[t]{0.53\linewidth}
\vspace*{-0.10in}
\raggedright
\tzpencyclopediaimage{1077}
\end{minipage}
\end{samepage}

\clearpage

\tzpentryheader{1077}{Direct Equation Progression and Proof Trajectory}

\tzpsection{Direct Equation Progression}
\begin{enumerate}[leftmargin=1.6em,itemsep=9pt,topsep=6pt]
\item \textbf{Aggregate Relation}
\[
M = \sum_{\mathrm{i}} \lvert i \rangle \langle i| \otimes \lvert i \rangle_{cl} \implies \text{Holevo Readout}
\]
This fixes the first explicit relation carried by the entry rather than a generic certificate record normalization.

\item \textbf{Support Relation}
\[
mathcal{J}_{\mathrm{out}} = \mathcal{T}_{\mathrm{TZP}}\!\left(\mathcal{K}_{\mathrm{TZP}}\right)
\]
This adds the next operational equation that advances the entry beyond its opening definition.

\item \textbf{Closure Relation}
\[
mathcal{A}_{\mathrm{adm}}\!\left(\mathcal{J}_{\mathrm{out}}\right)=1
\]
This closes the progression with an actual admissibility or closure relation instead of recycling the normalization shell.
\end{enumerate}

\tzpsection{Proof}
\textbf{Proof route.} show that the stated equations form a coherent progression for the entry and remain compatible under the TRAWIN ordering.

\textbf{Formal method.} prove that the final closure relation follows from the preceding entry-specific equations.

\medskip
\tzpproofartifacts{Isabelle audit: corpus classification recorded in appendix.}{Lean audit: compiled corpus certificate.}{Rocq audit: compiled corpus certificate.}

\tzpsection{Dependencies and Test Handle}
\begin{tzpsnapshot}
\tzpsnaprow{Dependencies}{\tzpidref{1076}, \tzpidref{1000}}
\tzpsnaprow{Node}{\tzpcode{registry_id_1077}}
\tzpsnaprow{Support Titles}{Aggregate Relation; Support Relation; Closure Relation}
\tzpsnaprow{Test Handle}{If the cited entry equations cannot be made mutually admissible in one TRAWIN-normalized frame, the entry fails.}
\end{tzpsnapshot}

\tzpsection{Canonical Equation Links}
\tzpequationlinks
  {K_{\mathrm{h}}=K_{L}}
  {\mathcal{J}_{\mathrm{out}} = \cdots}
  {\mathcal{A}_{\mathrm{adm}}\!\left(\mathcal{J}_{\mathrm{out}}\right) = \cdots}
